
\chapter{Идентификация клиентов и~ТСП}\label{ch:kyc-kyb}
Идентификация физлица (KYC, \textit{Know Your Customer}) и~идентификация юрлица
или ТСП (торгово-сервисное предприятие; KYB, \textit{Know Your Business})
формируют профиль клиента~--- совокупность собранных и~оценённых сведений,
на~которой строится решение по~каждой последующей операции. Комплаенс
в~платёжном бизнесе сводится к~трём вопросам: кого подключаем, что уже известно
о~клиенте или ТСП и~в~какой момент операцию нужно остановить.

\section{Идентификация при подключении физлиц}\label{sec:kyc-individuals}

\highlight{Идентификация физлица~--- это лестница уровней доверия}. На~каждой
ступени оператор решает, можно ли считать человека тем, за~кого он себя выдаёт,
и~какой риск принимает, открывая следующий уровень сервиса.

Лестница идёт от~минимального набора сведений к~документальной сверке
и~дополнительным проверкам там, где этого требует закон или профиль риска~---
набор признаков клиента, по~которым оператор относит его к~уровню риска,
заданному Положением Банка России от~15.10.2015~№~499-П
\enquote{Об~идентификации клиентов в~целях ПОД/ФТ}~\cite{cbr-499p}. Чем меньше
сведений собрано на~входе, тем уже набор доступных операций и~лимитов; Банк
России прямо разъясняет это на~примере переводов и~электронных кошельков по
упрощённой идентификации~\cite{cbr-banking-faq-identification}. Автоматизация
ускоряет вход, обязанность проверять качество данных остаётся: ошибка на~входе
редко ограничивается одним клиентом.

Стандартный набор проверок~--- сверка документов по~государственным источникам,
проверка по~перечням Росфинмониторинга и~иным санкционным спискам
(см.~гл.~\ref{ch:sanctions}), а~также выявление политически значимого лица
(PEP, \textit{politically exposed person})~--- человека, наделённого или ранее
наделённого значимой публичной функцией; в~процедурах противодействия
легализации (отмыванию) доходов и~финансированию терроризма (ПОД/ФТ) отдельно
учитывают также членов его семьи и~близких соратников и~применяют к~этой
категории усиленную проверку из-за повышенного коррупционного
риска~\cite{cbr-499p,fatf-pep-guidance}. Сама методология PEP восходит
к~рекомендациям FATF; о~приостановке членства РФ в~FATF и~текущей роли ЕАГ как
региональной \textit{FATF-style body} для СНГ~---
см.~раздел~\ref{sec:aml-fatf-eag}.

Для неидентифицированного пользователя электронных денежных средств (ЭДС)
161-ФЗ (статья~10) ограничивает остаток ЭДС суммой 15\,000~руб. в~любой момент
и~месячный оборот~--- суммой 40\,000~руб., и~эти ограничения действуют
до~прохождения упрощённой идентификации (без подтверждения личности
по~документу с~фотографией). После упрощённой идентификации лимит остатка ЭДС
поднимается до~100\,000~руб., а~месячный оборот~---
до~200\,000~руб.~\cite{fz-161}. Для полной (персонифицированной, то~есть
с~очной или эквивалентной сверкой документа, удостоверяющего личность)
идентификации физлица-резидента~РФ лимит остатка установлен законом на~уровне
600\,000~руб.; дополнительные параметры могут определяться договором
с~оператором. Продукт обязан хранить текущий статус идентификации клиента
и~применять соответствующие лимиты перед каждой операцией.

\subsection*{ЕСИА: упрощённая идентификация через инфраструктуру электронного правительства}

Основной канал упрощённой идентификации в~РФ~--- Единая система идентификации
и~аутентификации (ЕСИА), федеральная информационная система, созданная
Постановлением Правительства РФ от~28.11.2011~№~977 \enquote{О~ЕСИА} как часть
инфраструктуры электронного правительства~\cite{esia-pp977}. 115-ФЗ допускает
её использование кредитными организациями для упрощённой идентификации
физических лиц, ссылаясь на~инфраструктуру электронного правительства как
на~источник проверенных сведений~\cite{fz-115}.

Учётная запись ЕСИА имеет три уровня. \textbf{Упрощённая} (уровень~1)
подтверждена телефоном или email и~не~идентифицирует пользователя для целей
115-ФЗ. \textbf{Стандартная} (уровень~2) дополнительно требует СНИЛС
и~паспортные данные, прошедшие сверку с~реестрами Социального фонда России
и~МВД,~--- этого хватает для упрощённой идентификации в~смысле статьи~7 115-ФЗ.
\textbf{Подтверждённая} (уровень~3) требует личного визита в~центр обслуживания
или подтверждения кодом из~заказного письма Почты России и~открывает полный
объём государственных услуг.

В~инженерной практике PSP или банк получает результат проверки через систему
межведомственного электронного взаимодействия (СМЭВ): запрос возвращает набор
подтверждённых атрибутов клиента из~ЕСИА вместо очного предъявления документов.
Делегированная идентификация (Tinkoff~ID, Сбер~ID и~подобные) использует ту же
инфраструктуру: оператор-партнёр, уже идентифицировавший клиента, передаёт
подтверждённые атрибуты получающему банку с~согласия пользователя.

\subsection*{ЕБС: дистанционная полная идентификация по~биометрии}

Полная (персонифицированная) идентификация без~визита в~отделение банка
с~2018~года реализуется через Единую биометрическую систему (ЕБС). Нормативная
основа~--- Федеральный закон от~31.12.2017~№~482-ФЗ, изменивший 115-ФЗ и~закон
об~информации и~информационных технологиях~\cite{fz-482,fz-115}. Биометрия
(голос и~изображение лица) собирается очно в~отделении банка и~хранится
в~государственной системе. При следующем удалённом контакте клиент проходит
подтверждение по~биометрии и~получает услугу без~визита в~отделение.

Оператор ЕБС с~16~декабря 2022~года~--- АО \enquote{Центр биометрических
технологий}, созданное Указом Президента РФ от~30.09.2022~№~693; ранее функции
оператора исполнял ПАО \enquote{Ростелеком}~\cite{ebs-overview}. К~30~сентября
2023~года банки передали в~ЕБС накопленные массивы биометрических данных;
с~1~марта 2026~года микрофинансовые организации обязаны идентифицировать
заёмщиков по~ЕБС для дистанционной выдачи займов.

ЕБС открывает банку право предоставлять клиенту услуги, требующие полной
идентификации (повышение лимита остатка ЭДС до~600\,000~руб., открытие счёта,
выдача карты), без~очного визита. На~фоне зарубежных аналогов~--- индийского
Aadhaar, эстонского e-Estonia, бельгийского eID~--- ЕБС выделяется
централизованной государственной моделью с~банковской регистрацией: первое
предъявление биометрии всё равно требует физического присутствия
в~банке-партнёре.

\subsection*{Атаки на~биометрию: \textit{presentation} и~\textit{injection}}

\highlight{Требование первого очного визита в~отделение~--- контрмера к~атаке
вбросом (\textit{injection})}. \textit{Presentation attack}~--- предъявление
устройству захвата фальшивого артефакта: распечатки лица, экрана с~фотографией,
маски, видео на~стороннем дисплее. \textit{Injection attack}~--- обход
устройства захвата: атакующий вбрасывает синтетическое видео (дипфейк)
или~подменённый поток прямо в~API верификации, минуя камеру.

Стандарты охватывают эти классы по-разному. ISO/IEC~30107-3 описывает
методологию тестирования \textit{presentation attack detection} (PAD); iBeta
Level~1 и~Level~2~--- стандартная отраслевая сертификация для~PAD на~основе
ISO~30107-3~\cite{iso-30107-3}. \textit{Injection attacks} этим стандартом
не~покрыты. Европейская техническая спецификация CEN/TS~18099 от~2024~года
и~готовящийся ISO~25456 закрывают пробел~\cite{cen-ts-18099}; в~США NIST
в~четвёртой редакции Digital Identity Guidelines (SP~800-63-4),
финализированной в~августе 2025~года, формально вводит \textit{injection
attack} как отдельный класс~\cite{nist-sp-800-63-4-2025}.

Масштаб атак заметно вырос с~распространением AI-генерации лиц. Один из~крупных
финансовых институтов зафиксировал за~январь~--- август 2025~года~8065 попыток
обхода \textit{liveness} через дипфейки~--- порядок десятков попыток в~день
на~единичный канал~\cite{roc-injection-detection-2025}. Открытые платформы
\textit{deepfake-as-a-service} снизили порог входа до~любительского уровня;
статичная проверка живости (пассивный \textit{liveness}~--- по~артефактам
видеопотока) и~даже challenge-response (активный \textit{liveness}~---
моргнуть, повернуть голову) обходятся синтетическими моделями, работающими
в~реальном времени.

Архитектурный ответ~--- многофакторная (\textit{multi-factor}) верификация
на~стороне канала идентификации. ЕБС эту защиту реализует через первое очное
предъявление: атакующий должен скомпрометировать как биометрическую проверку,
так и~систему контроля доступа в~отделение банка. Чисто удалённое подключение
без~физического первого шага открывает класс атак вбросом, на~который
PAD-сертификация по~iBeta Level~2 ответа не~даёт. Устойчивость биометрического
канала в~CMS оценивается по~CEN/TS~18099 или~NIST~SP~800-63-4, не~по~устаревшей
ISO~30107-3.

\section{{EUDI} {Wallet} и~{eIDAS}~2.0}\label{sec:kyc-eudi-wallet}

\highlight{К~сентябрю 2026~года все страны~ЕС обязаны выдавать гражданам
и~резидентам European Digital Identity Wallet}
(EUDI~Wallet)~\cite{eu-eidas-2-regulation}. Это часть Регламента 2024/1183,
известного как eIDAS~2.0, заменившего регламент~910/2014 в~части цифровой
идентичности. С~2027~года регулируемые отрасли~--- банки, телеком, госуслуги,
страхование~--- обязаны принимать кошелёк как канал аутентификации.

EUDI~Wallet работает по~модели избирательного раскрытия (\textit{selective
disclosure}): пользователь сам решает, какие атрибуты раскрыть проверяющей
стороне (\textit{relying party}). Можно подтвердить факт совершеннолетия
без~раскрытия даты рождения, факт резидентства без~раскрытия адреса, факт
юридической дееспособности без~раскрытия профессии. Технически это реализуется
через подтверждаемые удостоверения (\textit{verifiable credentials})
с~криптографическими доказательствами с~нулевым раскрытием знания (ZKP-классы).
Архитектурная и~эталонная модель (ARF)~--- техническая спецификация,
на~2025~год действует версия~2.8~--- описывает протоколы и~форматы
для~эмитентов, кошельков и~проверяющих сторон~\cite{eudi-wallet-arf}.

ЕБС~--- централизованная государственная система с~единым реестром
биометрических шаблонов; EUDI~Wallet~--- децентрализованная модель, где кошелёк
находится на~устройстве пользователя, а~эмитент атрибута (государство, банк,
университет) подписывает credential, не~получая информации о~местах его
использования. ЕБС оптимизируется под~единый канал верификации; EUDI~Wallet~---
под~интероперабельность между независимыми сервисами и~странами.

Для~российского PSP с~европейскими клиентами или~корреспондентскими отношениями
EUDI~Wallet с~2027~года становится вторым стандартом удалённой идентификации
параллельно с~ЕБС. Системы хранения профиля клиента в~CMS должны поддерживать
оба формата: атрибуты из~ЕБС по~существующему формату и~подтверждаемые
удостоверения из~EUDI~Wallet по~ARF-схеме.

\section{Идентификация при подключении ТСП}\label{sec:kyb-merchants}

У~юридического лица проверяют и~личность, и~деловую конструкцию: структуру
собственности, реальный вид деятельности, возвраты, спорные операции и~риск
использования бизнеса как транзитной оболочки. \highlight{Идентификация юрлица
  собирает регистрационные данные, владельцев и~платёжное поведение в~единый
профиль риска}.

Базовый пакет документов включает выписку из~реестра, данные директора
и~учредителей, сведения о~бенефициарном владельце (\textit{ultimate beneficial
owner}, UBO), банковские реквизиты и~описание бизнеса. Нестыковки в~структуре
владения или сайт, не~похожий на~заявленный вид деятельности,~--- повод для
дополнительных вопросов. Документы у~рискованного бизнеса часто формально
в~порядке: после подключения профиль сверяют с~живым поведением (наблюдаемыми
транзакционными паттернами). Логика идентификации по~115-ФЗ~\cite{fz-115}
и~499-П это прямо предусматривает: \highlight{кредитная организация должна
  установить данные самого клиента и~сведения о~бенефициарном
владельце}~\cite{cbr-499p}.

Код категории торгового предприятия (\textit{Merchant Category Code}, MCC) сам
по~себе нейтрален; правила схем требуют корректно классифицировать ТСП
и~передавать правильный код категории. Для категорий повышенного риска
(\textit{high-integrity risk merchants} в~терминологии Visa) действуют
отдельные ограничения и~требования мониторинга~\cite{visa-core-rules}. Подмена
MCC~--- попытка спрятать реальный профиль бизнеса под~более мягкой категорией.
Если заявленная категория не~сходится с~фактической деятельностью, платёжный
поток останавливают до~выяснения.

Модель платёжного фасилитатора (\textit{payment facilitator}, PayFac) переносит
ответственность за~идентификацию субмерчантов (\textit{sponsored merchants}~---
продавцов, подключённых через PayFac), их~KYB, текущий мониторинг и~разбор
споров на~саму платформу; содержание проверки остаётся прежним. Visa возлагает
на~эквайера ответственность за~PayFac и~его субмерчантов: до~начала обработки
нужно подтвердить регистрацию PayFac у~Visa, присвоить уникальные
идентификаторы самому PayFac и~его субмерчантам, получать отчётность по~их
активности и~следить за~корректным MCC~\cite{visa-core-rules}. Платформа
подключает субмерчантов быстрее, расплачивается собственными процедурами
идентификации и~более плотным мониторингом. Без них быстрое подключение
превращается в~быстрое накопление риска.

Глубина должной проверки (\textit{due diligence})~--- комплекса процедур
по~сбору и~оценке сведений о~клиенте~--- растёт вместе с~профилем риска ТСП:
расширяется пакет документов, учащается мониторинг, к~высокорисковым категориям
применяют плавающий резерв (\textit{rolling reserve})~--- удержание части
оборота ТСП на~отдельном счёте до~закрытия окна спора.

\practicenote{%
  При проверке ТСП главный маркер~--- несоответствие заявленного MCC
  фактической деятельности.
}

Проверка ТСП собирается из~трёх шагов: вход (документы, UBO, банковские
реквизиты в~единый профиль), решение (риск-ориентированный подход, RBA~---
  подробнее см.~гл.~\ref{ch:aml}~--- присваивает уровень риска и~определяет
глубину должной проверки) и~постпроверка (мониторинг сравнивает заявленное
с~наблюдаемым). Все три шага опираются на~один объект данных~--- запись ТСП
в~CMS, где регистрационные сведения, бенефициары, MCC, лимиты и~резерв связаны
общим идентификатором. Несоответствие MCC заявленной деятельности система ловит
автоматически только при~наличии второго независимого источника данных
о~фактической специализации продавца: онбординг эквайера, страница оплаты,
поведение чеков. Без~такого источника MCC остаётся самозаявленным, а~сверка
переходит к~ручному разбору после первой жалобы. PayFac-модель добавляет
к~записи ещё~один уровень: атрибуты субмерчантов хранятся в~CMS платформы
на~той~же глубине, что и~прямых ТСП эквайера, иначе единый идентификатор
перестаёт работать как точка сверки.

\needspace{9\baselineskip}
\begin{table}[H]
  \begingroup
  \scriptsize
  \setlength{\tabcolsep}{4pt}
  \renewcommand{\arraystretch}{1.2}
  \begin{minipage}{\textwidth}
    \begin{tabularx}{\textwidth}{
        @{}B{0.18\textwidth}
        Y
        Y
        Y@{}
      }
      \toprule
      \headrow
      \thd{Критерий} &
      \thd{Низкий риск} &
      \thd{Средний риск} &
      \thd{Высокий риск} \\
      \midrule
      Примеры MCC &
      retail, SaaS (Software as a Service) &
      travel &
      crypto exchange,\newline gambling, adult \\
      Документы &
      стандартный пакет &
      расширенный пакет\newline + лицензии &
      расширенный пакет\newline + финансовая отчётность \\
      Мониторинг после подключения &
      квартальный &
      ежемесячный &
      еженедельный \\
      Плавающий резерв &
      нет &
      5--10\% &
      10--20\% \\
      \bottomrule
    \end{tabularx}
  \end{minipage}%
  \endgroup
  \caption{Список для идентификации ТСП по~категориям риска.}\label{tab:kyb-risk}
\end{table}

\section{Мониторинг ТСП после подключения}\label{sec:merchant-monitoring}

Риск меняется после подключения, и~мониторинг продолжает отслеживать изменения
объёма, географии и~доли возвратов. Схемы типа bust-out (когда клиент или~ТСП
  длительное время ведёт себя добросовестно, чтобы нарастить лимиты, а~затем
резко уводит средства; подробнее~--- в~гл.~\ref{ch:fraud}) и~попытки
использовать ТСП для~отмывания проявляются именно здесь. Пока профиль выглядит
спокойно, мониторинг ослабляет внимание, а~проблема становится видна поздно.

Поведенческие признаки для расследования~--- резкий рост оборота, нетипичная
география клиентов, высокий процент возвратов и~несоответствие MCC реальной
деятельности.

В~промышленных системах это надстройка над~авторизационной базой: агрегаты
по~ТСП, скользящие окна, поиск выбросов и~очередь случаев для аналитиков.
Резкий рост объёма у~ТСП команда видит почти сразу, а~не в~конце месяца.
В~европейском регионе Visa требует от~эквайера сохранять ежедневные данные
по~обороту, среднему чеку, количеству операций, времени до~расчёта и~количеству
диспутов и~сравнивать их с~нормальным профилем активности минимум
ежедневно~\cite{visa-core-rules}.

\subsection*{Цикл обновления данных по~уровню риска}

\highlight{499-П задаёт периодичность обновления сведений о~клиенте по~уровню
риска}: для клиентов высокого и~среднего риска~--- не~реже раза в~год, для
клиентов низкого риска~--- не~реже раза в~три года~\cite{cbr-499p}.
Внеочередное обновление обязательно при изменении сведений о~клиенте и~при
возникновении подозрений в~отношении его операций. Аналогичная логика
применяется и~к~ТСП: плановый пересмотр KYB-профиля по~уровню риска плюс
внеочередной триггер по~событию.

В~ПОД/ФТ-периметре цикл обновления~--- самый частый пропуск. Идентификация
при~подключении проходит через единый шлюз с~интеграциями (ЕСИА, ЕБС, проверка
по~реестрам, формирование профиля). Обновление по~графику живёт в~отдельной
очереди задач, куда клиент попадает с~приближением даты по~таймеру. Очередь
распределяется по~аналитикам KYC/KYB-команды и~при высокой нагрузке
откладывается~--- особенно для среднего риска, где запас в~год выглядит
достаточным до~конца квартала, а~к~декабрю обновление превращается в~массовую
кампанию с~резким падением качества разбора.

Инженерное решение: на~каждом клиентском профиле в~CMS поведение определяют две
даты~--- последнего обновления и~контрольного срока следующего. Триггер
поднимает задачу за~две-четыре недели до~контрольной даты. Без таймера в~записи
клиента обновление превращается в~годовой отчёт, а~годовой отчёт превращается
в~предписание регулятора.

\section{Платформа \enquote{Знай своего клиента}}\label{sec:kyc-zsk}

\highlight{Банк России с~июля 2022~года ведёт собственную платформу оценки
риска юридических лиц и~ИП~--- \enquote{Знай своего клиента}, ЗСК}. Каждый
клиент получает одну из~трёх зон риска: зелёная~--- без~ограничений, жёлтая~---
усиленное наблюдение, красная~--- жёсткие ограничения вплоть до~отказа
в~обслуживании и~расторжения договора банковского
счёта~\cite{cbr-platform-zsk}. Оценка строится на~совокупности признаков:
профиль операций, контрагенты, отраслевая принадлежность, история регулятивных
обращений, связи с~уже маркированными лицами.

С~1~октября 2024~года Банк России открыл бесплатный публичный доступ к~сервису
проверки контрагентов~\cite{cbr-zsk-public-2024}. До~этой даты доступ имели
только кредитные организации; теперь любая компания может запросить категорию
ЗСК потенциального контрагента, и~это смещает часть KYB-нагрузки на~ТСП.
Контрагент в~красной зоне~--- маркер высокого риска для~продавца независимо
от~внутреннего KYB; контрагент в~жёлтой зоне требует расширенной проверки
документов и~внеочередного обновления профиля.

Инженерно ЗСК~--- внешний источник данных для~внутреннего RBA банка, а~не
замена KYB. Категория ЗСК поступает асинхронно от~Банка России, оператор
сравнивает её с~собственной оценкой клиента. Расхождения разбираются вручную:
если банк присвоил клиенту низкий риск, а~ЗСК вывел в~красную зону, оператор
обязан пересмотреть профиль и~провести дополнительные процедуры по~115-ФЗ.
Обратный сценарий (клиент в~зелёной зоне ЗСК, но~банк нашёл признаки высокого
риска) права на~снижение не~даёт: внутренний RBA остаётся ведущим.

В~феврале 2025~года Банк России анонсировал аналог ЗСК для~физических лиц,
направленный на~снижение числа ошибочных блокировок счетов добросовестных
клиентов. Архитектурно это та же модель: платформа с~централизованным реестром
оценок риска, доступом для~кредитных организаций и~асинхронным потоком
обновления категорий. ЗСК для~физлиц войдёт в~ту же интеграционную точку CMS,
что и~текущая ЗСК для~юрлиц, но~по~отдельной таблице с~собственной
классификацией.

\section{KYC/KYB как готовый сервис}\label{sec:kyc-industry-providers}

Большинство команд не~пишет KYC-конвейер с~нуля. Сегмент
\textit{compliance-as-a-service} закрывает идентификацию физлиц и~юрлиц через
единый API: верификация документа с~селфи-фотофиксацией, биометрическое
сравнение, проверка по~санкционным и~PEP-спискам, сверка по~реестрам
юридических лиц. Под одной интеграцией оказываются десятки источников данных,
которые иначе пришлось бы подключать поочерёдно.

Из крупных игроков рынка~--- Sumsub~\cite{sumsub-overview}, зарегистрированный
в~Лондоне в~2015~году выходцами из~Санкт-Петербурга под маркой
Sum~\&~Substance: ранние клиенты приходили из~СНГ, сейчас платформа обслуживает
банки и~финтехи в~десятках юрисдикций.
ComplyAdvantage~\cite{complyadvantage-overview}~--- лондонская платформа
финмониторинга: санкционный скрининг (\textit{sanctions screening}),
мониторинг транзакций (\textit{transaction monitoring}) и~поиск негативных
упоминаний (\textit{adverse media}) в~одном API. На~западных рынках в~том же
классе работают Jumio, Onfido, Persona~--- с~акцентом на~верификацию документов
и~биометрию.

Выбор \enquote{строить или купить} определяется объёмом потока, профилем риска
клиентского сегмента и~ценой ошибки. Готовый сервис закрывает 80--90~\% типовых
сценариев и~заметно ускоряет регуляторное прохождение; собственный конвейер
оправдан там, где специфика клиентской воронки и~рисков не~покрывается
стандартными правилами провайдера~--- у~крупных банков с~собственным аппетитом
к~риску и~у~финтехов, чей основной продукт~--- собственно идентификация.

\section{{Travel} {Rule} для {VASP}}\label{sec:kyc-travel-rule}

\highlight{Recommendation~16 FATF (Travel Rule) переносит требование KYC
из~банковского периметра в~криптообмены}. VASP (\textit{Virtual Asset Service
Provider}~--- криптобиржа, кастодиальный кошелёк, OTC-деск) обязан при~переводе
криптоактивов между двумя VASP передавать получающей стороне набор атрибутов
отправителя и~получателя: имя, адрес, идентификатор кошелька, иногда дата
рождения. Данные передаются вне~блокчейна (\textit{off-chain}) параллельно
с~ончейн-транзакцией через протоколы вроде TRP или~Sumsub Travel Rule
и~используются получающим VASP для~санкционного скрининга и~записи в~аудит.

На~2025~год Recommendation~16 имплементировали~85 юрисдикций, ещё~14 находятся
в~процессе~\cite{fatf-targeted-update-va-2025}. Порог применения различается:
в~ЕС и~Великобритании~--- 0~евро (применяется к~любой сумме), в~США~--- \$3000,
в~большинстве других юрисдикций~--- эквивалент~\$1000. В~ЕС действует
с~30~декабря 2024~года EU \textit{Transfer of Funds Regulation} (TFR,
регламент~2023/1113) вместе с~финальными руководящими указаниями
(\textit{guidelines}) EBA от~июля 2024~года~\cite{eba-travel-rule-2024}.

Для~российского PSP, работающего с~криптоэквайрингом через зарубежных
VASP-партнёров, это становится дополнительным KYC-требованием на~стороне
партнёра. VASP-партнёр запрашивает у~PSP атрибуты плательщика по~Travel Rule
до~зачисления криптоактива; если данные неполны или~не~прошли санкционный
скрининг получающей стороны, перевод останавливается. Российская сторона пока
не~имплементировала Recommendation~16 во~внутреннее законодательство;
для~трансграничных операций значение имеют требования партнёра в~зарубежной
юрисдикции, не~российский норматив.
